\documentclass[border=10pt]{standalone}
\usepackage[no-math]{fontspec}
\usepackage[brazil]{babel}
\usepackage{amsmath,amssymb}
\IfFileExists{C:/Windows/Fonts/Lato_Regular.ttf}{%
  \setmainfont{Lato}[
    Path = C:/Windows/Fonts/,
    UprightFont = Lato_Regular.ttf,
    BoldFont = Lato_Bold.ttf,
    ItalicFont = Lato_Italic.ttf,
    BoldItalicFont = Lato_BoldItalic.ttf
  ]
  \setsansfont{Lato}[
    Path = C:/Windows/Fonts/,
    UprightFont = Lato_Regular.ttf,
    BoldFont = Lato_Bold.ttf,
    ItalicFont = Lato_Italic.ttf,
    BoldItalicFont = Lato_BoldItalic.ttf
  ]
}{%
  \RequirePackage[default]{lato}
}
\renewcommand{\familydefault}{\sfdefault}

\usepackage{xcolor}
\usepackage{tikz}
\usepackage{pgfplots}
\pgfplotsset{compat=1.18}
\usetikzlibrary{arrows.meta,positioning,shapes.geometric,fit,backgrounds,calc,decorations.pathreplacing,patterns}

% Paleta institucional UEFS / Bioengenharia de Solos
\definecolor{bsNavy}{HTML}{1F4268}
\definecolor{bsNavyDark}{HTML}{152E4A}
\definecolor{bsGold}{HTML}{B38A3D}
\definecolor{bsGoldDark}{HTML}{8C6927}
\definecolor{bsTeal}{HTML}{0E7C7B}
\definecolor{bsForest}{HTML}{2D6A4F}
\definecolor{bsAlert}{HTML}{9E2A2B}
\definecolor{bsEarth}{HTML}{8C5A3C}
\definecolor{bsWater}{HTML}{2E86AB}
\definecolor{bsInk}{HTML}{26313D}
\definecolor{bsMuted}{HTML}{5B6573}

\definecolor{bgNavy}{HTML}{EDF2F7}
\definecolor{bgGold}{HTML}{F8F3E9}
\definecolor{bgTeal}{HTML}{EBF6F5}
\definecolor{bgForest}{HTML}{EEF5F0}
\definecolor{bgAlert}{HTML}{FDF2F2}

\tikzset{
  >=Stealth,
  bsBox/.style={
    draw=#1, fill=white, rounded corners=5pt, line width=1.15pt,
    inner sep=7pt, align=center, font=\footnotesize\sffamily, text=bsInk
  },
  bsTint/.style={
    draw=#1, fill=#1!8!white, rounded corners=5pt, line width=1.15pt,
    inner sep=7pt, align=center, font=\footnotesize\sffamily, text=bsInk
  },
  bsBadge/.style={
    fill=#1, text=white, rounded corners=9pt, inner xsep=8pt, inner ysep=3pt,
    font=\scriptsize\bfseries\sffamily
  },
  bsTitle/.style={font=\normalsize\bfseries\sffamily, text=bsNavy, align=center},
  bsSub/.style={font=\footnotesize\sffamily, text=bsMuted, align=center},
  bsArrow/.style={->, >=Stealth, line width=1.35pt},
  bsDArrow/.style={<->, >=Stealth, line width=1.35pt}
}


\begin{document}
\begin{tikzpicture}

  % Parâmetros de escala do Gantt
  % Coluna de rótulos: x = -7.7 a 0.0
  % Meses M1..M12: cada mês = 0.92 cm -> x = 0.0 a 11.04
  \def\colw{0.92}
  \def\rowh{0.78}

  % Título
  \node[bsTitle] at (1.67, 5.85)
    {Cronograma Típico de Implantação Integrada de Bioengenharia (Meses 1 a 12)};
  \node[bsSub] at (1.67, 5.40)
    {Sequenciamento construtivo: obras hidráulicas e contenções na estiagem, revegetação no período úmido};

  % Fundo do cabeçalho dos meses
  \fill[bsNavy] (0, 4.25) rectangle (11.04, 4.85);
  \fill[bsNavyDark] (-7.7, 4.25) rectangle (0, 4.85);
  \node[font=\scriptsize\bfseries, text=white] at (-3.85, 4.55) {FASE / FRENTE DE SERVIÇO};

  \foreach \m in {1,...,12} {
    \pgfmathsetmacro{\xc}{(\m - 0.5)*\colw}
    \node[font=\scriptsize\bfseries, text=white] at (\xc, 4.55) {M\m};
  }

  % Fundo alternado das 6 linhas + grade vertical
  \foreach \r in {1,...,6} {
    \pgfmathsetmacro{\ytop}{4.25 - (\r - 1)*\rowh}
    \pgfmathsetmacro{\ybot}{4.25 - \r*\rowh}
    \ifodd\r
      \fill[bgNavy!55!white] (-7.7, \ybot) rectangle (11.04, \ytop);
    \else
      \fill[white] (-7.7, \ybot) rectangle (11.04, \ytop);
    \fi
    \draw[bsMuted!25!white, line width=0.5pt] (-7.7, \ybot) -- (11.04, \ybot);
  }

  \foreach \m in {0,...,12} {
    \pgfmathsetmacro{\xg}{\m*\colw}
    \draw[bsMuted!28!white, line width=0.5pt] (\xg, -0.43) -- (\xg, 4.25);
  }

  % Borda externa da tabela
  \draw[bsNavy!60!white, line width=1pt] (-7.7, -0.43) rectangle (11.04, 4.85);
  \draw[bsNavy!60!white, line width=0.9pt] (0, -0.43) -- (0, 4.85);

  % Helper macro para linha de Gantt
  \newcommand{\ganttrow}[6]{
    \pgfmathsetmacro{\ymid}{4.25 - (#1 - 0.5)*\rowh}
    \pgfmathsetmacro{\xstart}{(#3 - 1)*\colw + 0.06}
    \pgfmathsetmacro{\xend}{#4*\colw - 0.06}
    \node[anchor=west, font=\scriptsize\bfseries, text=bsInk] at (-7.50, \ymid) {#2};
    \fill[#5, rounded corners=3.5pt] (\xstart, \ymid - 0.21) rectangle (\xend, \ymid + 0.21);
    \node[font=\tiny\bfseries, text=white] at ({(\xstart+\xend)/2}, \ymid) {#6};
  }

  \ganttrow{1}{1. Levantamento, Topografia \& Projeto}{1}{2}{bsNavy}{M1--M2}
  \ganttrow{2}{2. Drenagem de Crista \& Canaleta Verde}{2}{4}{bsWater}{M2--M4}
  \ganttrow{3}{3. Estruturas de Base (Krainer / Enrocamento)}{3}{5}{bsEarth}{M3--M5}
  \ganttrow{4}{4. Obras de Vertente (Paliçadas / Fascinas / Biomantas)}{4}{6}{bsTeal}{M4--M6}
  \ganttrow{5}{5. Semeadura, Estacas Vivas \& Plantio Ciliar}{5}{8}{bsForest}{M5--M8}
  \ganttrow{6}{6. Manutenção \& Monitoramento Pós-Chuva}{7}{12}{bsGold}{M7--M12}

  % Faixas sazonais na base (largura com folga para o texto)
  \fill[bgGold, draw=bsGold, rounded corners=3pt, line width=0.8pt] (0.04, -1.15) rectangle (4.54, -0.58);
  \node[font=\scriptsize\bfseries, text=bsGoldDark, align=center] at (2.29, -0.865) {Estação Seca:\\[-1pt]Obras \& Contenções};

  \fill[bgForest, draw=bsForest, rounded corners=3pt, line width=0.8pt] (4.66, -1.15) rectangle (11.00, -0.58);
  \node[font=\scriptsize\bfseries, text=bsForest, align=center] at (7.83, -0.865) {Estação Chuvosa:\\[-1pt]Estabelecimento Vegetal \& Monitoramento};

\end{tikzpicture}
\end{document}
